\subsection{Simultaneous perfect quotients and a further mixed extension}

\begin{statement}{D.40}{Proved}{An extension through every finite perfect source quotient}
Let $\mathscr P$ be the set of finite-index normal subgroups
$M$ of the marked group $B=C_2*C_3$ for which $B/M$ is
perfect. Define the actual marked quotient
\[
 \mathcal R=\overline{\operatorname{im}\!
     \left(B\longrightarrow\prod_{M\in\mathscr P}B/M\right)},
 \qquad
 K_{\mathcal R}=\bigcap_{M\in\mathscr P}\overline M
       \ \subset\widehat B.                              \tag{D.40.1}
\]
Then $\mathcal R=\widehat B/K_{\mathcal R}$ and its maximal
prosolvable quotient is exactly
\[
 \mathcal R/\mathcal R'\simeq B^{\rm ab}\simeq C_6.        \tag{D.40.2}
\]
Here $\mathcal R'$ denotes the closed subgroup generated
by commutators. Its finite images are perfect.
The isomorphism in (D.40.2) uses the original
marked images of $s$ and $r$, of orders two and three.

Write $R_T=\overline{[L,L]}\cap\overline{D_9K}$ for D.39's
source kernel and retain its normalized automorphism
$\alpha_{39}$ of $\mathcal T=\widehat B/R_T$. There are
canonical marked isomorphisms
\[
 \begin{split}
 \mathcal U
   &=\widehat B/(R_T\cap K_{\mathcal R})
     \simeq\mathcal T\mathbin{\times_{C_6}}\mathcal R,\\
 \mathcal V
   &=\widehat B/(R_T\cap K_{\mathcal R}\cap\ker\sigma_7)
     \simeq\mathcal U\mathbin{\times_{C_2}}S_7,
 \end{split}                                             \tag{D.40.3}
\]
where $\sigma_7$ is the actual surjection
\[
 s\longmapsto(1\,2)(3\,4)(5\,6),\qquad
 r\longmapsto(1\,4\,6)(2\,7\,5),                          \tag{D.40.4}
\]
and the last fibre product uses the common sign map.
On these groups the maps
\[
 (\alpha_{39},1_{\mathcal R}),\qquad
 (\alpha_{39},1_{\mathcal R},1_{S_7})                    \tag{D.40.5}
\]
are actual continuous peripheral automorphisms with
multiplier one and the two hexagon equations. Their
simultaneous loop conjugators, common gauge and
commutator are constructed below. Thus the D.39 action
extends coherently through every finite perfect quotient
of the full marked source and through (D.40.4).

In particular it acts as the identity on
\[
 \mathcal P_{\ge5}
    =\prod_{p\ge5}\operatorname{PSL}_2(\mathbb Z_p)        \tag{D.40.6}
\]
under the actual modular maps
$s\mapsto\left[\begin{smallmatrix}0&1\\-1&0\end{smallmatrix}\right]$,
$r\mapsto\left[\begin{smallmatrix}0&-1\\1&-1\end{smallmatrix}\right]$.
The resulting projection
$\mathcal V\to\mathcal T\times\mathcal P_{\ge5}$ is onto.
Its source kernel in $\widehat B$ before passing to
$\mathcal V$ is exactly $R_T\cap\ker\rho_{\ge5}$.
Some additional exact projection kernels are
\[
 \begin{split}
 \ker(\mathcal U\to\mathcal T)&=\mathcal R',\\
 \ker(\mathcal V\to\mathcal U)&=A_7,\\
 \ker(\mathcal V\to\mathcal T)&=\mathcal R'\times A_7.
 \end{split}                                             \tag{D.40.7}
\]

The remaining full-source kernel
$R_T\cap K_{\mathcal R}\cap\ker\sigma_7$ is nontrivial.
Its restriction under the actual surjection
\[
 \begin{split}
 \sigma_{10}(s)&=(1\,2)(3\,4)(5\,6)(7\,8)(9\,10),\\
 \sigma_{10}(r)&=(1\,3\,9)(4\,10\,8)(5\,6\,7)
 \end{split}                                             \tag{D.40.8}
\]
maps onto $A_{10}$. This is a computed remaining source
kernel, not a nonlifting assertion for the complete CM
pair. D.45 excludes a full profinite lift of this recorded
action; a replacement realizing the same CM pair remains open.
The original $\widehat F$ source maps onto
$\ker(\mathcal V\to S_3)$ with the same displayed kernel.
The compact-cover and fixed-fibre CM data remain exactly
those of D.39.

This addresses Part D(e), section 51.16, folio 734 verso,
scan 174, and D-after-LM2(1), section 51.15, folios
725--727, scans 156--160. The finite-subdirect argument,
the two explicit central covers, and all source maps
are proved here.
\end{statement}
\begin{proof}
\emph{The finite-subdirect argument.}
Let $H$ be the image of $B$ in a finite product
$G_1\times\cdots\times G_j$ of its perfect quotients.
Every coordinate map is onto. Write $H_\infty$ for
the stable term of the finite derived series. It is
perfect and $H/H_\infty$ is the maximal solvable quotient.
For each coordinate kernel $N_i$, the image of
$H_\infty$ is $G_i$: every term of the derived series
has that image, and the series stabilizes. Hence
$H=H_\infty N_i$, so every $N_i$ maps onto
$H/H_\infty$. On the other hand
\[
 [N_1,N_2,\ldots,N_j]\subseteq\bigcap_iN_i=1.
                                                               \tag{D.40.9}
\]
For $j=1$ this just says that the single coordinate
kernel is trivial. Taking images in $H/H_\infty$ shows
that this finite solvable quotient is nilpotent.

A finite nilpotent quotient of $B$ is generated by
an element of order dividing two and an element of
order dividing three. They lie in different Sylow
factors and commute. The quotient is therefore cyclic
of order dividing six. In particular $H'=H_\infty$
is perfect. Every finite quotient of $\mathcal R$
factors through one such finite coordinate image:
an open kernel in the closed product image contains
the intersection of finitely many coordinate kernels.
It follows that every finite solvable quotient of
$\mathcal R$ is a quotient of $C_6$, and that every
finite image of $\mathcal R'$ is perfect.

\emph{Determining both factors of the common cyclic quotient.}
The upper bound six is attained. The following two
marked matrix pairs give actual surjections from $B$:
\[
 \begin{array}{c|c|c}
 G&s&r\\ \hline
 \operatorname{SL}_3(\mathbb F_7)&
 \begin{pmatrix}6&0&0\\0&6&0\\0&0&1\end{pmatrix}&
 \begin{pmatrix}4&5&1\\1&2&3\\2&0&1\end{pmatrix}\\[7pt]
 \operatorname{SL}_4(\mathbb F_3)&
 \begin{pmatrix}2&0&0&0\\0&2&0&0\\0&0&1&0\\0&0&0&1\end{pmatrix}&
 \begin{pmatrix}0&2&2&1\\0&0&1&0\\1&0&0&2\\0&1&2&1\end{pmatrix}
 \end{array}                                             \tag{D.40.10}
\]
Both pairs have determinant one and satisfy $s^2=r^3=1$.
Their generation is verified by actual words and orbit
lower bounds. In their faithful actions on nonzero
vectors, the recorded stabilizer-chain orbit lengths are
\[
 \begin{array}{c|c|r}
 G&\text{orbit lengths}&|G|\\ \hline
 \operatorname{SL}_3(\mathbb F_7)&342,\ 336,\ 49&5\,630\,688\\
 \operatorname{SL}_4(\mathbb F_3)&80,\ 78,\ 72,\ 27&12\,130\,560 .
 \end{array}                                             \tag{D.40.11}
\]
Each generator used in these orbit calculations is
stored as a word in the displayed $s,r$. A separate
standard-library verifier evaluates all those words
as matrices, checks their permutations, and computes
the orbits of the generators fixing all earlier base
points. Each orbit is a lower bound for the orbit
of the actual stabilizer in $\langle s,r\rangle$.
The products in (D.40.11) equal the upper bounds
\[
 |\operatorname{SL}_n(\mathbb F_q)|
   =q^{n(n-1)/2}\prod_{a=2}^n(q^a-1),
\]
obtained by counting ordered bases and dividing by
the $q-1$ determinant values. Thus the displayed
words generate the full groups; no bare group-order
output is needed for this inference.

For all matrix commutator identities in this proof use
$[a,b]=aba^{-1}b^{-1}$. These two groups are perfect
by elementary matrix identities. Row elimination, using elementary
transvections and
\[
 \begin{pmatrix}a&0\\0&a^{-1}\end{pmatrix}
 =
 \left(E_{12}(a)E_{21}(-a^{-1})E_{12}(a)\right)
 \left(E_{12}(-1)E_{21}(1)E_{12}(-1)\right),
\]
shows that elementary matrices generate each special
linear group. For three distinct indices,
$[E_{ij}(a),E_{jk}(b)]=E_{ik}(ab)$.
Every elementary generator in dimensions three and
four is therefore a commutator. The verifier also
checks these identities over both fields.

In the first group, $z=2I$ is central of order three;
replace the displayed $r$ by $zr$. In the second,
$z=-I$ is central of order two; replace $s$ by $zs$.
The torsion relations still hold. Each new map
$\phi_1$ is a central-character twist of the original
epimorphism $\phi_0$. On $B'$ the two maps agree,
and $\phi_0(B')=G$ because $G$ is perfect. Thus the
twisted maps are also onto. Their generation is
additionally verified by the same word-and-orbit
procedure.

For each pair the joint image is $G\times C_d$,
$d=3$ or $2$, under
$(g,\epsilon)\mapsto(g,z^\epsilon g)$. Indeed
$B\to G\times C_d$ is onto: a common quotient
of the perfect group $G$ and the cyclic group is
trivial, so the elementary subdirect-product argument
gives the full product. Equivalently, the ratio
$\phi_1(b)\phi_0(b)^{-1}$ gives the specified central
character. On $s,r$ the two characters are exactly
\[
 \chi_2(s)=1,\quad\chi_2(r)=0,\qquad
 \chi_3(s)=0,\quad\chi_3(r)=1.                            \tag{D.40.12}
\]
All four perfect quotients are coordinates of
$\mathcal R$. Their central ratios therefore
give a surjection $\mathcal R\to C_2\times C_3=C_6$.
Together with the upper bound already proved, this
establishes (D.40.2), with its exact marked map.

\emph{The common quotient and simultaneous conjugators.}
Every finite quotient of $\mathcal T$ is solvable.
The group $H$ is an extension of $Q$ by
$C_2^2\times C_3$, and $Q$ is solvable: reduction
modulo four has two-group kernel over $S_3$,
while $\operatorname{SL}_2(\mathbb F_3)$ has
normal $Q_8$ with cyclic quotient. D.39's
homology group is therefore an extension of a
solvable finite group by an abelian profinite
group. Also $T_3$ is an extension of the
solvable $Q_0$ by a finite two-group.
The displayed embedding of $\mathcal T$ in
their product proves the assertion.
Consequently every common marked quotient
of $\mathcal T$ and $\mathcal R$ factors through
the just-computed $C_6$. Conversely $R_T$ is
contained in $\widehat F'\subset\widehat B'$, so
$\mathcal T$ retains precisely the original map
$B\to B^{\rm ab}=C_6$. The maximal common marked
quotient is therefore exactly $C_6$.

For completeness, the subdirect-product argument
used here has explicit kernels. For a closed
subgroup $J\subset X\times Y$ surjecting onto both
factors, set
$K_X=\{x:(x,1)\in J\}$ and
$K_Y=\{y:(1,y)\in J\}$. These are closed normal
subgroups and the pairs in $J$ induce a canonical
isomorphism $X/K_X\simeq Y/K_Y$. Its inverse
identifies $J$ with that fibre product. Applying
this to the image of $\widehat B$ proves the first
isomorphism of (D.40.3), including the quotient
by the intersection of the two source kernels.

The automorphism $\alpha_{39}$ fixes $B^{\rm ab}$:
its images of $s,r$ are their conjugates. It hence
pairs with $1_{\mathcal R}$ on the common $C_6$.
There are also compatible actual conjugators, not
only compatible automorphisms. Let $A,\beta$ be
the raw D.39 conjugators and let
$k=\beta_y^{\rm ab}-\beta_x^{\rm ab}\in\widehat{\mathbb Z}$
be its actual global gauge. Their retained $H$
coordinates are $(6,0)$, and $k\equiv0\pmod6$.
Since $x=t^{-2}$ and $y=sxs$ have images $(0,1)$
in $C_2\times C_3=B^{\rm ab}$, an element of $F$
with abelianization $(a,b)$ maps to $(0,a+b\bmod3)$.
Thus
\[
 A,\ \beta,\ x^kA,\ x^k\beta
       \quad\hbox{all have image zero in }B^{\rm ab}.
                                                               \tag{D.40.13}
\]
The finite seed checks these residues, and D.39's
integral and odd-prime limits retain them.

Before gauge, use the paired conjugators
\[
 (A,x_{\mathcal R}^{-k}),\qquad
 (\beta,x_{\mathcal R}^{-k}).                            \tag{D.40.14}
\]
Their two components agree in $C_6$, by (D.40.13).
On $\mathcal R$ they give the inner automorphism
$\operatorname{Int}(x_{\mathcal R}^{-k})$, which
fixes the cusp. The common gauge by $x^k$ changes
them to $(x^kA,1)$ and $(x^k\beta,1)$, and changes
the automorphism to $(\alpha_{39},1)$.
The resulting $f$ is $(f_{39},1)$.

Both conjugators lie in the image of $\widehat F$,
as seen under the retained map to $S_3$.
The full source kernel is contained in $\widehat F'$,
so the marked abelianization remains
$\widehat{\mathbb Z}^2$. D.39's gauge calculation
then gives an actual lift of this common $f$ in
$\widehat F'$. The torsion and exact cusp equations
give the same two $m=0$ hexagons. The component
automorphisms and their inverses preserve the
fibre product, proving generation in every finite
coordinate image and compatibility in the limit.

\emph{Every good-prime congruence level.}
Use the exact two matrices in (D.40.6). Their
product is $\left(\begin{smallmatrix}1&-1\\0&1\end{smallmatrix}\right)$,
so $x$ is $\left(\begin{smallmatrix}1&2\\0&1\end{smallmatrix}\right)$.
These are the same modular markings as D.25.
They generate every projective special linear
quotient modulo an odd integer: their elementary
upper and lower unipotents generate the local
matrix groups by row elimination, and the Chinese
remainder map supplies the independent local
coefficients.

For $p\ge5$ and every $e\ge1$, put
$D=\operatorname{diag}(2,2^{-1})$ over $\mathbb Z/p^e$.
The exact commutator identities
\[
 [D,E_{12}(a/3)]=E_{12}(a),\qquad
 [D,E_{21}(-4a/3)]=E_{21}(a)                            \tag{D.40.15}
\]
show that $\operatorname{SL}_2(\mathbb Z/p^e)$,
and its projective quotient, are perfect. Their
canonical kernels therefore belong to $\mathscr P$.
Projection to these coordinates and passage to
the inverse limit give the actual surjection
$\mathcal R\to\mathcal P_{\ge5}$. Surjectivity
onto the whole product follows from the same
Chinese remainder generation at every finite
collection of primes and exponents.

Every finite quotient of $\mathcal P_{\ge5}$ is
perfect, while every finite quotient of
$\mathcal T$ is solvable. Their only common
quotient is trivial, since a finite perfect
solvable group is trivial and profinite groups
are separated by finite quotients. The same
subdirect-product argument proves that their
actual joint source image is
$\mathcal T\times\mathcal P_{\ge5}$, with source
kernel $R_T\cap\ker\rho_{\ge5}$. Formula (D.40.5)
acts as the identity on the congruence component.
The proof uses the actual global $k$: its odd
components are $-6$. It does not substitute the
integer finite-seed gauge $-438$ at unrelated
prime levels.

\emph{The $S_7$ extension and the next explicit kernel.}
The certificate enumerates all $5040$ elements
generated by (D.40.4), and all $3\,628\,800$
elements generated by (D.40.8). The torsion
relations are checked independently. In the
first quotient $t$ has cycle type $(4,3)$ and
$x=t^{-2}$ has order six. In the second $t$
has cycle type $(4,5,1)$ and $x$ has order ten.

We give the elementary normal-subgroup facts
needed for the kernel argument. For $n\ge5$,
every nontrivial normal subgroup of $A_n$
contains a three-cycle. A permutation with a
cycle of length at least four gives a three-cycle
by taking its commutator with a three-cycle on
three consecutive entries. For a permutation
having a three-cycle and other entries, the
commutator with $(a\,b\,d)$, where $a,b$ are
successive entries and $d$ is outside that cycle,
gives either a five-cycle or a double transposition.
A permutation consisting of transpositions gives
a double transposition by the same calculation
on two of its pairs. A double transposition,
commuted with a three-cycle using one entry of
each pair and a fifth entry, gives a five-cycle,
reducing to the first case. A sole three-cycle
already has the required form.

All three-cycles are conjugate in $A_n$: correct
an odd conjugating permutation by a transposition
of two unused entries. They generate $A_n$, since
each pair of transpositions is a product of
three-cycles. Thus $A_n$ is simple and perfect.
An element centralizing $A_n$ in $S_n$ fixes every
three-cycle and hence every three-element support;
intersecting such supports fixes every point.
Its centralizer is trivial. It follows that the
normal subgroups of $S_n$ are exactly
$1,A_n,S_n$: a normal subgroup intersecting
$A_n$ trivially centralizes it.

There is no continuous surjection $\mathcal U\to S_7$.
Such a map would factor through a finite marked
image
$J\subset T_0\times G_1\times\cdots\times G_j$,
where $T_0$ is a finite solvable quotient of
$\mathcal T$ and each $G_i$ is perfect. Let
$N_0,N_1,\ldots,N_j$ be the coordinate kernels.
The image of $N_0$ in $S_7$ contains $A_7$,
because its quotient is solvable. For $i\ge1$,
$J'$ maps onto $G_i$, so $J=J'N_i$.
The normal image of $N_i$ therefore contains
an odd permutation and must be all of $S_7$.
But the iterated commutator of these coordinate
kernels lies in their trivial intersection,
whereas its image contains $A_7$. This is a
contradiction.

The largest common marked quotient of $\mathcal U$
and the actual $S_7$ quotient is consequently
$C_2$: the only alternatives among quotients
of $S_7$ are $S_7$ and the trivial group, and
the common sign quotient is retained. This
proves the second isomorphism in (D.40.3).
The action $\alpha_{39}$ fixes that sign, so
pairing with the identity gives the second
automorphism in (D.40.5). The normalized
conjugators in (D.40.13) pair with the identity
conjugators in $S_7$. Before gauge the $S_7$
component of $x^{-k}$ is already the identity,
because $x$ has order six and $k\equiv0\pmod6$.
This proves the simultaneous marking, the
gauge and the hexagons in this further extension.
Taking the kernels of its explicit fibre-product
projections gives (D.40.7).

Finally there is no surjection $\mathcal V\to S_{10}$.
Factor a proposed map through a finite image
\[
 J\subset T_0\times G_1\times\cdots\times G_j\times S_7.
                                                               \tag{D.40.16}
\]
The $T_0$-coordinate kernel has image containing
$A_{10}$. The $S_7$-coordinate kernel also has
image containing $A_{10}$: otherwise it would
be trivial, forcing a surjection from $S_7$
onto the larger group $S_{10}$.
Every perfect-coordinate kernel has image
$S_{10}$, by the same $J=J'N_i$ argument.
The iterated commutator of all these kernels
is trivial in (D.40.16), but maps onto the
nontrivial perfect group $A_{10}$. This
contradiction proves the claim.

Under $\sigma_{10}$ the kernel
$R_T\cap K_{\mathcal R}\cap\ker\sigma_7$ has
normal image in $S_{10}$. It is contained in
$\widehat F'\subset\widehat B'$, so this
image is contained in $A_{10}$. The image
cannot be trivial, since that would make
$\sigma_{10}$ factor through $\mathcal V$.
Simplicity therefore makes the image exactly
$A_{10}$. The exact source kernel and the
open embedding of $\widehat F$ give the final
source map in the statement. All previous CM
evaluations factor through the unchanged
$\mathcal T$ component. The nonzero remaining
kernel and its actual alternating quotient
identify a further full-source test; they do
not assert a failure of every possible full
lift of the CM pair.
\end{proof}
